\chapter{Sharing data and results safely}\label{ch:sharing}

A partner wants the analytic file, the lawyer wants it de-identified, and both want it by Friday. When you share data or publish a table, you take on a duty to the people the evaluation studied, and \enquote{we removed the names} is not nearly enough. In this chapter we examine disclosure risks and build checks to support a documented release decision.

We cover the quasi-identifiers that make removing names insufficient, how to suppress small cells without lying with the totals, how to generate synthetic stand-ins for the real records, and how to gate raw intake files so likely protected fields are identified before further processing. \marginnote{All four of these techniques, from the quasi-identifier scan to the intake gate, protect the people behind the data, and doing that well is what keeps a partner willing to share the next file.} Throughout, you measure rather than assert: you can compute useful indicators of disclosure risk, provided you state which records and quasi-identifiers they cover, and once it is a number a lawyer and an evaluator can agree on a threshold and check it. By the end you will be able to compute that number for a file you are about to share, coarsen and suppress until it clears the threshold a data-use agreement set, sweep a raw intake folder before anything imports from it, generate a synthetic stand-in when the real records cannot leave at all, and file the one-page disclosure record naming what shipped and under what standard. Do that and the file reaches the partner by Friday with the lawyer's signature on it, and a year later you can still say what left the building and under what rule. Figure~\ref{fig:ch14-viz-pipeline} distinguishes the checks for a record-level file from those for a published table.

\section{Quasi-identifiers: why removing names is not enough}\label{sec:qid}
Removing direct identifiers leaves a dataset far more exposed than it looks, because combinations of ordinary variables can single a person out. The classic trap is that ZIP code, birth date, and sex together identify a large share of the population, so a \enquote{de-identified} file with those three is not de-identified.\sidenote{\textcite{sweeney2002kanon} put a number on it: roughly 87\% of Americans are uniquely identifiable from five-digit ZIP, full date of birth, and sex alone. She re-identified the Massachusetts governor's health record from a \enquote{de-identified} release using nothing more.} We measure the exposure with $k$-anonymity\index{k-anonymity@$k$-anonymity} (how many people share each combination of quasi-identifiers\index{quasi-identifiers}) and $l$-diversity\index{l-diversity@$l$-diversity} (how varied the sensitive values are within each such group), and we suppress or coarsen until the risk is acceptable.\sidenote{The exposure is not hypothetical. \textcite{narayanan2008deanon} re-identified named individuals in the \enquote{anonymized} Netflix Prize ratings by matching a handful of public movie ratings against the release, showing that sparse, high-dimensional records stay linkable even after every direct identifier is stripped. For an evaluator whose caseload is exactly that shape, a few hundred people rated across many program touchpoints, the lesson is direct: the columns you kept for analysis are the columns an outsider can match on.} These metrics give an embedded evaluator a defensible, replicable standard that a practitioner and a lawyer can both check, and health-informatics teams have already run them as a working de-identification method on real releases, not just a theoretical ideal \parencite{elemam2008kanon}.

\begin{figure*}[t]
\centering
\begin{tikzpicture}[font=\footnotesize]
\node[io] (file) at (0,0) {the analytic\\file};
\node[stage] (scan) at (3.1,1.1) {\texttt{riskscan}:\\$k$ and $l$};
\node[decision] (dk) at (6.0,1.1) {$k \ge$\\threshold?};
\node[stage] (co) at (6.0,3.0) {coarsen\\(\S\ref{sec:coarsen})};
\node[stage] (tab) at (3.1,-1.2) {tabulate for\\publication};
\node[stage] (sup) at (6.4,-1.2) {\texttt{suppress}: primary\\+ complementary};
\node[stage] (chk) at (9.9,0) {disclosure-review\\checklist (six fields)};
\node[io] (rel) at (13.0,0) {the release};
\draw[flow] (file) -- (scan);
\draw[flow] (scan) -- (dk);
\draw[flow] (dk) -- node[right=1pt] {\scriptsize no} (co);
\draw[flow] (co) -| (scan);
\draw[flow] (dk) -- node[above=2pt, xshift=3mm] {\scriptsize yes} (chk);
\draw[flow] (file) -- (tab);
\draw[flow] (tab) -- (sup);
\draw[flow] (sup) -- (chk);
\draw[flow] (chk) -- (rel);
\node[substep, text width=4.8cm, anchor=north] (raw) at (1.8,-2.5) {upstream of everything, \texttt{rawsweep} (\S\ref{sec:rawsweep}) sweeps the raw intake, to flag likely protected fields for review before import};
\draw[flow, dashed] (raw.north) -- (file.south);
\end{tikzpicture}
\caption[The release pipeline this chapter teaches]{The chapter as one pipeline. A file being shared takes the top lane: scan with \texttt{riskscan}, coarsen and rescan until the smallest quasi-identifier cell clears the agreed threshold. A table being published takes the bottom lane: blank the small cells, plus the complementary blanks that keep printed totals from undoing them. Both products stop at the same one-page checklist before anything leaves, and a synthetic stand-in (\S\ref{sec:synth}) requires its own disclosure review when real records cannot be shared.}
\label{fig:ch14-viz-pipeline}
\end{figure*}


The scan is only as good as the list it runs on, so settle the quasi-identifier list with the partner before anyone computes anything. You can screen a candidate column with a blunt test: could a coworker, a neighbor, or a public record supply this value for a specific person? Age, ZIP, race, job, and program site usually pass that test; a survey attitude item usually does not. Write the agreed list down, with the date and who agreed to it, because it becomes the first line of the disclosure-review checklist this chapter assembles and the first thing a lawyer asks to see.

Each metric is a single number the release script can compute. A file satisfies $k$-anonymity at threshold $k_{\min}$ when the \emph{smallest} quasi-identifier group is still at least that big,
\[
  k \;=\; \min_{g \in G}\, \bigl|\{\, i : q_i = g \,\}\bigr| \;\ge\; k_{\min},
\]
where $G$ is the set of distinct quasi-identifier combinations, $q_i$ is record $i$'s quasi-identifier tuple, and $\lvert\cdot\rvert$ counts the records sharing a combination; a cell with $k=1$ is one person standing alone. The $l$-diversity check adds a second condition, that every group contain at least $l$ distinct sensitive values,
\[
  \min_{g \in G}\, \bigl|\{\, s_i : q_i = g \,\}\bigr| \;\ge\; l,
\]
where $s_i$ is record $i$'s sensitive value\marginnote{A group can be large in $k$ and still leak: if every member shares one diagnosis ($l=1$), then knowing a person is in the group is knowing their diagnosis.}.

You can run both checks with one packaged command, and the box below gives its install line, its syntax, and the results it returns.

\begin{kaobox}[frametitle=Package Integration: \texttt{riskscan}]
\textbf{Install and run:} \texttt{riskscan}\index{riskscan@\texttt{riskscan}} scans a dataset for the quasi-identifiers you name, computes $k$-anonymity across their combinations, and flags the records at highest re-identification risk, so the exposure is measured rather than assumed away. Install it with the house idiom of Appendix~\ref{app:toolkit} (repository \texttt{riskscan-stata-public}).
The syntax is \texttt{riskscan varlist [if] [in], [k(\#) flag(newvar) sensitive(varname) detail]}. It returns \texttt{r(cells)}, \texttt{r(k1)}, and \texttt{r(below)} (the count under the threshold), so a release script can branch on the number rather than a human's read of a table. Pass \texttt{sensitive()} to add the $l$-diversity check; pass \texttt{detail} for the highest-risk cells (the listing caps at 30, and it never prints the sensitive values themselves).
\end{kaobox}

An agreed threshold makes the scan reproducible, but passing it is only part of disclosure review. Check the proposed fields, possible outside information, sensitive attributes, and earlier releases before deciding whether the file can be shared.\marginnote{HIPAA accepts two routes to a de-identified file. Safe Harbor requires removal of specified identifiers and no actual knowledge that the remaining information could identify an individual. Expert Determination requires an appropriately qualified expert to assess and document a very small identification risk. A chosen $k$ threshold alone does not establish either route. Knowing which standard your data-use agreement invokes tells you whether a checklist or a computation is what you owe.} \index{re-identification}

\subsection{A worked example: how many people are one of a kind?}\label{sec:kanon}
A data steward about to release a file wants a blunt answer before signing off: once the names are stripped, how many records are unique on the columns that remain, and so potentially linkable by someone who knows those facts and can establish that the person is in the file? We answer it on \texttt{sysuse nlsw88}\index{sysuse@\texttt{sysuse}}, the 1988 wage survey that ships with Stata, treated here as a stand-in for an administrative caseload of 2{,}246 records. Nothing in it is secret, which is exactly why it is safe to demonstrate on; the arithmetic is identical when the records are real clients. The familiarity is also the lesson: these are the same 2{,}246 women whose wages we decomposed in Chapter~\ref{ch:claims} and charted in Chapter~\ref{ch:graphs}, and the columns that made them analyzable all book long are exactly the columns that make them findable here. We pick four quasi-identifiers that any partner file would plausibly contain (race, marital status, college-graduate status, and industry) and count how many records share each combination:\marginnote{A larger cell gives a record more peers on the fields being assessed. It does not rule out identification from other fields, outside information, or knowledge that a person belongs in the file.}
\par
{\footnotesize
\begin{verbatim}
sysuse nlsw88, clear
egen cell = group(race married collgrad ///
    industry), missing
bysort cell: gen k = _N
label var k "people sharing this quasi-id cell"
\end{verbatim}
}
The variable \texttt{cell} is one number per distinct combination of the four columns, and \texttt{k} is how many people fall in that combination. A person with \texttt{k}${}=1$ is unique on those four columns: removing their name does not remove that uniqueness, although successful identification also depends on outside information and knowledge of membership\marginnote{Race, marital status, degree, and industry are four facts a coworker or a neighbor could supply without ever seeing the file.}. We bin the records by cell size and count both cells and people in each bin:
\par
{\footnotesize
\begin{verbatim}
gen kbin = 1 + (k>1) + (k>4) + (k>10)
label define kb 1 "k=1 (unique)" 2 "k=2-4" ///
    3 "k=5-10" 4 "k>10", replace
label values kbin kb
tabulate kbin            // people per bin
preserve
    bysort cell: keep if _n==1
    tabulate kbin        // cells per bin
restore
\end{verbatim}
}
The first \texttt{tabulate} counts people because each row is a person. After \texttt{bysort cell: keep if \_n==1}, each row is a distinct combination, so the second table counts cells. That \texttt{keep} changes the data in memory. Preserve the person-level file before this display and restore it before calling \texttt{riskscan} or computing the coarsened counts below; otherwise you will scan one representative per cell and obtain the wrong result. Table~\ref{tab:kanon} reports both: the number of distinct quasi-identifier cells in each size band, and the number of people those cells contain.

\begin{table*}[htbp]
\centering
{\footnotesize
\begin{tabular}{lrr}
\hline
Cell size ($k$) & Cells & People \\
\hline
$k=1$ (unique)   &  24 &    24 \\
$k=2$--$4$        &  25 &    63 \\
$k=5$--$10$       &  11 &    74 \\
$k>10$            &  43 & 2{,}085 \\
\hline
Total            & 103 & 2{,}246 \\
\hline
\end{tabular}
}
\caption[$k$-anonymity of four ordinary columns]{Re-identification exposure of \texttt{nlsw88} (an administrative stand-in, $N=2{,}246$) on four quasi-identifiers: race, married, collgrad, industry. The four columns carve the file into just 103 distinct cells. Twenty-four people are unique on these four columns alone, and another 63 sit in cells of four or fewer, so 87 of the 2{,}246 people, spread across 49 of the 103 cells, sit in thin cells. Built by \texttt{ch14\_kanon.do}.}
\label{tab:kanon}
\end{table*}

The count that matters is the first row, and it is worth printing on its own:
\par
{\footnotesize
\begin{verbatim}
count if k==1
\end{verbatim}
}
\par
{\footnotesize
\begin{verbatim}
  24
\end{verbatim}
}
Running the same scan as a packaged command reproduces the table above exactly, so you can trust that the tool and the hand-computation agree:
\par
{\footnotesize
\begin{verbatim}
riskscan race married collgrad industry
\end{verbatim}
}
prints the same four bins, reports 103 distinct cells and 24 unique records, and stores 87, the count of people in cells thinner than five, in \texttt{r(below)}. Because the count is stored in a return scalar rather than only printed on screen, a release script can gate on it without a human reading a table: \texttt{if r(below) > 0} then hold the release. That is the practical use of turning within-file uniqueness into a number \parencite{elemam2008kanon}: the release decision runs as code.

The scan finds twenty-four records unique on race, marital status, college graduation, and industry. This is evidence of uniqueness within the released file, rather than a demonstration that an outsider can identify each person. Another 63 sit in cells of two to four; an additional field such as age or city could further divide those groups. \marginnote{The arithmetic: the four columns have $3\times2\times2\times12$ possible combinations, but only 103 are actually occupied, so about 22 people share the average occupied cell. The distribution is lumpy, and the rare combinations (a college graduate in a small industry, say) collapse to one.} Plot those cell sizes and the imbalance is obvious: dozens of one- and two-record cells cluster at the left of Figure~\ref{fig:kdist}, while a few large cells on the right contain most of the file. The risk assessment should concentrate on that thin left tail and on the outside information available to the intended recipient.

Translate the result for a non-technical review: \enquote{Twenty-four records have unique combinations of these four fields within this file. Whether someone could name the people depends on outside information, knowledge of file membership, and a successful linkage.} Record the fields, file version, intended recipient, and script with that statement.

\begin{figure*}[htbp]
\centering
\par
{\footnotesize
\begin{verbatim}
bysort cell: keep if _n==1
gen kplot = min(k, 30)   // pool cells of 30+ at the right
histogram kplot, discrete frequency ///
    xtitle("Cell size k")
\end{verbatim}
}
\includegraphics[width=\linewidth]{ch14_kdist.png}
\caption[Distribution of quasi-identifier cell sizes]{How many cells have each size, over race, married, collgrad, and industry in the \texttt{nlsw88} stand-in. The tall spike at the left is the danger: 24 cells contain a single person and 15 more contain just two, so dozens of records are uniquely or nearly-uniquely identified by four ordinary columns. Cells of 30 or more people are pooled at the right, where these quasi-identifiers distinguish fewer people. The mass on the left is what coarsening and suppression exist to move rightward. Built by \texttt{ch14\_kanon.do}.}
\label{fig:kdist}
\end{figure*}

Before anyone signs, go back through the assumptions behind the number: the box below names the three conditions and the situation in which each one fails.

\begin{restson}[What the $k$-anonymity number depends on]
\begin{enumerate}
\item The attacker holds only the quasi-identifiers you counted. \emph{Noticed by:} a partner file arriving with columns your scan never named.
\item $k$ is counted on the released file, not the population. \emph{Noticed by:} a subset release, where a cell of five in your full file can be a cell of one in the extract.
\item Every release of the same records is counted together. \emph{Noticed by:} two safe-looking cuts that intersect to a cell of one.
\end{enumerate}
\end{restson}

\subsection{Coarsening: the cheapest way to raise \texorpdfstring{$k$}{k}}\label{sec:coarsen}
When you need to empty the $k=1$ bin, the quickest move you can make is to coarsen the mesh, because a record is unique only relative to the resolution of its columns. One \texttt{recode} does most of the work in this section, and by the end you will be able to state the trade in both directions: how many one-of-a-kind people a coarser column protects, and which comparison the reader gives up in exchange.\marginnote{Coarsen the highest-cardinality column first: the singletons cluster in the rare combinations, and those almost always involve the variable with the most categories. Age in single years, five-digit ZIP, and detailed occupation codes are the usual culprits; bin them before you reach for anything more elaborate.} Industry is the finest of our four quasi-identifiers, twelve categories, so it does the most damage; collapsing it to three broad groups (goods-producing, services, and public administration) merges many singleton cells into populated ones without touching the other three columns:
\par
{\footnotesize
\begin{verbatim}
recode industry (1/4=1 "Goods") ///
    (5/11=2 "Services") (12=3 "Public"), ///
    gen(ind3)
egen cell3 = group(race married collgrad ind3), ///
    missing
bysort cell3: gen k3 = _N
count if k3==1
\end{verbatim}
}
Coarsening\index{coarsening} one column moves people off the singleton spike wholesale. Table~\ref{tab:coarsen} shows the before-and-after: the same four quasi-identifiers, but with industry at twelve categories and then at three.

\begin{table*}[htbp]
\centering
{\footnotesize
\begin{tabular}{lrr}
\hline
                       & $k=1$ people & Distinct cells \\
\hline
Industry, 12 groups    & 24 & 103 \\
Industry, 3 groups     &  3 &  39 \\
\hline
\end{tabular}
}
\caption[Coarsening industry collapses unique records]{Collapsing industry from twelve categories to three, holding race, married, and collgrad fixed, cuts the number of one-of-a-kind people from 24 to 3 and the number of distinct cells from 103 to 39. The other three columns are untouched; the whole gain comes from one coarser variable. Built by \texttt{ch14\_kanon.do}.}
\label{tab:coarsen}
\end{table*}

Collapse industry from twelve groups to three and the unique records fall from 24 to 3, an eight-fold reduction, at the cost of some analytic resolution. \marginnote{The mechanism is plain: a coarser column is a wider net, so more people share each combination, so fewer stand alone.} The trade is explicit and it is the analyst's to make: three industry groups may be plenty for a report that only ever compares goods to services, in which case the coarsening costs nothing the reader would have used. When it does cost something, coarsening and suppression combine, coarsen until the remaining singletons are few, then suppress those. The value here is that the trade is visible and defensible: you can tell a lawyer exactly how many people the release protects and exactly what analytic detail that protection cost.

Coarsening is also the first of three exits we lay out in Section~\ref{sec:supexample}: coarsen the columns, suppress the cells, or aggregate the whole table to a higher level. When you have to choose among them, start by asking the person who will read the release: which comparison does the report need at full resolution? Get the answer in writing before the first \texttt{recode} runs.

\section{Suppressing small cells without lying}\label{sec:suppress}
Small cells create substantial disclosure risk, and blanking one cell may fail when row and column totals let a reader recover its value. Use primary suppression for the small cell, then complementary suppression where visible totals would permit reconstruction.\marginnote{Complementary suppression and $l$-diversity address different risks: the former limits reconstruction of hidden counts from totals; the latter examines variation in sensitive values within quasi-identifier groups.} You can pass the marked output column that results into \texttt{collect}\index{collect@\texttt{collect}}, Stata's built-in table-assembly system (Stata~17 and later), or into another export path.

The threshold itself is not the analyst's to invent: most data-use agreements name one (five and ten are common), and the release record should cite the document that supplied it. Translate the $l$-diversity finding the same way you translated the $k$ count in Section~\ref{sec:kanon}, into a sentence a non-technical partner can act on. A cell where $l=1$ means \enquote{everyone in this group shares the same diagnosis, so knowing someone is in the group is knowing their diagnosis}, a sentence more persuasive than the inequality it summarizes.

Rates require two separate checks. For statistical stability, apply a denominator rule set for the intended use and show the denominator with the rate. For confidentiality, test whether the published count, denominator, rounded rate, and totals allow a reader to reconstruct a value below the release threshold. For example, \texttt{replace rate = . if denom < 30} implements a chosen stability rule; it does not by itself establish that the table is safe to release.

The box below introduces \texttt{suppress}, the packaged command that automates both kinds of blanking.

\begin{kaobox}[frametitle=Package Integration: \texttt{suppress}]
\textbf{Install and run:} \texttt{suppress}\index{suppress@\texttt{suppress}} automates public-release QA for a table: primary small-cell suppression, complementary suppression so totals do not reveal the hidden cells, and a marked output column an analyst can publish. Install it with the house idiom of Appendix~\ref{app:toolkit} (repository \texttt{suppress-stata-public}).
The syntax is \texttt{suppress countvar [if] [in], threshold(\#) [by(varlist) gens(newvar) complementary]}. It protects one grouping dimension per run, so a two-way table published with both row and column totals needs one run in each direction with the union of the flags blanked. It returns \texttt{r(n\_primary)} and \texttt{r(n\_complementary)}, and \texttt{gens()} writes a display column that prints the primary blanks as \texttt{<\#} and the complementary blanks as \texttt{*}, so the two kinds of suppression stay legible in the released table rather than collapsing to indistinguishable holes.
\end{kaobox}

\subsection{A worked example: blanking a cell without leaking it}\label{sec:supexample}
The problem is easiest to see in a two-way table small enough to hold in your head, so we blank one by hand before running the packaged command on it, and you will finish able to publish a table whose blanks cannot be differenced back and whose totals still add up. Staying in the stand-in file, cross college-graduate status against the four goods-producing industries, where the thin cells cluster, and count people. Suppose the release rule is the common one: blank any cell with fewer than five people. Here is the raw count:
\par
{\footnotesize
\begin{verbatim}
tabulate collgrad industry if industry<=4
\end{verbatim}
}
\par
{\footnotesize
\begin{verbatim}
          |            Industry
  College |    Ag  Mining  Constr  Manuf |  Total
----------+------------------------------+-------
 not grad |    14       4      26    334 |    378
     grad |     3       0       3     33 |     39
----------+------------------------------+-------
    Total |    17       4      29    367 |    417
\end{verbatim}
}
Under the rule implemented by \texttt{suppress}, positive counts below five are primary suppressions: three college graduates in agriculture, three in construction, and four non-graduates in mining. Zero is excluded from this primary rule; the zero in mining will be hidden as a complementary suppression. Decide explicitly whether your release policy treats zeros differently. But blanking them alone leaks them, because the reader has the column totals. The agriculture column totals 17 and its other cell is 14, so the hidden graduate cell is exactly $17-14=3$, recovered by subtraction; construction leaks the same way ($29-26=3$). \marginnote{Suppressing a cell while leaving the arithmetic that reconstructs it is not suppression at all.}

You close that leak with \emph{complementary} suppression\index{complementary suppression}: in any column that has just one blank left visible against its total, blank the surviving cell too, so the total can no longer be differenced back. Agriculture and construction each have one primary blank and one visible cell, so the algorithm blanks the visible non-graduate cell in both (14 and 26). Mining also needs its zero graduate cell hidden because the other cell is a primary suppression. Table~\ref{tab:suppress} shows the published version.

\begin{table*}[htbp]
\centering
{\footnotesize
\begin{tabular}{lrrrrr}
\hline
College   & Ag & Mining & Constr & Manuf & Total \\
\hline
not grad  & \textbf{--} & \textbf{--} & \textbf{--} & 334 & 378 \\
grad      & \textbf{--} & \textbf{--} & \textbf{--} &  33 &  39 \\
\hline
Total     & 17 &  4 & 29 & 367 & 417 \\
\hline
\end{tabular}
}
\caption[Primary and complementary suppression]{The same cross-tabulation after suppression. The primary blanks are the three positive cells below five (grad/Ag, grad/Construction, and notgrad/Mining). The complementary blanks are notgrad/Ag, notgrad/Construction, and the zero grad/Mining cell: without them, each column total would reveal the graduate cell by subtraction. Only manufacturing, where both cells are comfortably large, prints numbers. The totals are published unchanged, so the table still adds up honestly. Built by \texttt{ch14\_kanon.do}.}
\label{tab:suppress}
\end{table*}

Now every column with a small cell has both of its cells hidden, so no total can be differenced to a single value\marginnote{In the published table the agriculture column shows only its total of 17 split between two blanks, so a reader cannot tell whether the graduate cell is 1, 2, or 3.}. The hidden positive counts cannot be recovered by subtracting the other cell in their column from its total. This local check does not establish safety against all combinations of published margins or outside information. A single blank in a line of visible numbers is no secret: the total minus the numbers you can see is the blank. Hiding one small cell therefore means hiding a second, so the reader is left with two unknowns and one equation. The totals stay truthful, and that is the standard that matters: a suppressed table that still adds up is publishable, and one that quietly alters totals to hide a cell has crossed from suppression into fabrication.

Before you publish any table that prints totals, watch out for one particular trap.

\begin{vizcallout}[Suppression that leaks: the single-blank trap]
The most common disclosure error in a published table is a single suppressed cell in a fully-totaled row or column. If every other cell in the line is visible and the total is visible, the blank is not hidden, it is defined: it equals the total minus the visible cells. Any release with row and column totals needs at least two blanks in every line that has one, and a quick audit catches the leak before it ships: for each suppressed cell, confirm that a second suppression exists in both its row and its column. The \texttt{suppress} program checks the groups supplied in \texttt{by()}; it does not automatically inspect every other table or margin you may release.
\end{vizcallout}

Inside \texttt{suppress}, complementary suppression runs as a loop: over each \texttt{by()} group, while the group still contains exactly one suppressed cell, the program blanks the next-smallest surviving cell until at least two are hidden. Figure~\ref{fig:ch14-viz-suppress} traces that decision for a single cell.

\begin{figure*}[htbp]
\centering
\begin{tikzpicture}[node distance=10mm and 16mm]
\node[io] (cell) {each cell};
\node[decision, right=of cell] (d1) {cell \\ $0 < n < t$?};
\node[stage, right=24mm of d1] (prim) {primary\\suppress};
\node[decision, below=16mm of prim] (d2) {one blank\\left in\\its line?};
\node[stage, left=24mm of d2] (comp) {blank next-\\smallest cell};
\node[io, below=12mm of d2] (pub) {publish\\line};

\draw[flow] (cell) -- (d1);
\draw[flow] (d1) -- node[annot,above]{yes} (prim);
\draw[flow] (d1) |- node[annot,pos=0.22,left]{no} (pub);
\draw[flow] (prim) -- (d2);
\draw[flow] (d2) -- node[annot,above]{yes} (comp);
\draw[flow] (comp) -- ($(comp.north)+(0,6mm)$) -| node[annot,pos=0.25,above]{recheck} (d2.north);
\draw[flow] (d2) -- node[annot,right]{no} (pub);
\end{tikzpicture}
\caption[The suppression decision, cell by cell]{What the routine decides for each cell. A positive cell below the threshold $t$ is a primary suppression; the routine re-checks the line, and while it finds exactly \emph{one} blank against a visible total, it blanks the next-smallest surviving cell too, so the total can no longer be differenced back to the hidden value. A cell at or above $t$ may still be hidden as a complementary suppression. Every released line needs review against its visible totals. The \texttt{recheck} loop is why a very thin group can consume several complementary blanks before it is safe.}
\label{fig:ch14-viz-suppress}
\end{figure*}

Run against the goods-producing cross-tabulation, \texttt{suppress} reproduces the hand-computed table above cell by cell:
\par
{\footnotesize
\begin{verbatim}
* contract: one row per collgrad-industry cell, count in n
contract collgrad industry if industry<=4, ///
    zero freq(n)
suppress n, threshold(5) by(industry) ///
    complementary gens(n_pub)
\end{verbatim}
}
reports three primary suppressions and three complementary ones, marks agriculture, mining, and construction as protected, and leaves manufacturing clean, exactly the published version in Table~\ref{tab:suppress}. The zero cell (no college graduates in mining) is not itself a primary suppression, under the positive-count rule used by this command, but the routine can still draw it as complementary cover. One caveat applies to the method and the tool alike: if a group's blanked cells sum to exactly one, a reader deduces the split must be $\{1,0\}$, so a group that thin is not fully protected by suppression and needs coarsening or a higher threshold instead.

Figure~\ref{fig:ch14-viz-suppress} shows the mechanics; the judgment is choosing which exit to take when a table fails. In the goods-producing table, three moves were available. Suppression, the one taken, kept the four-industry breakdown but blanked six of eight cells, defensible if manufacturing provides the comparison the reader needs. Coarsening would have merged the four industries into one goods-producing column, publishing every count but surrendering the breakdown entirely. Aggregating would have kept the industry detail but reported it only at a higher level, statewide rather than by site, say. The tiebreaker is the reader's actual question: if the argument turns on construction versus manufacturing, six blanks gut the table and coarsening or a higher-level cut serves better; if manufacturing is the story, the suppressed version loses nothing of use. Whichever exit you take, log the cells affected and the reason; that entry is what the release record points to.

Suppress this way and you can publish the small-area dashboard policy audiences ask for without exposing the few people in a thin cell.\marginnote{Federal statistical agencies have moved past cell suppression toward differential privacy, which adds calibrated random noise to limit how much including or excluding one person changes the probability of possible released outputs, for exactly this reason: suppression protects against differencing one table, but a determined analyst who combines many overlapping released tables can still triangulate. Review even a single report for overlapping tables and external information; repeated releases add further opportunities to infer hidden values.} Access and privacy end up in the same table, and neither had to give way.

\subsection{Where the field is going: formal guarantees beyond \texorpdfstring{$k$}{k}-anonymity}\label{sec:frontier-dp}
Suppression and $k$-anonymity protect against differencing one table, but they make no promise once many overlapping releases pile up, which is where the field has moved. \textcite{dwork2006diffpriv} introduced differential privacy to close exactly that gap: rather than reasoning about which columns an adversary might hold, it adds calibrated random noise so that including or excluding one person changes the probabilities of possible outputs within a specified bound, with protection that accounts for outside information \parencite{dworkroth2014found}. \marginnote{The practical signal that differential privacy is not just theory: the U.S.\ Census Bureau adopted it for the 2020 census after internal research showed the old cell-suppression rules could be reverse-engineered to reconstruct individual records \parencite{abowd2018census}.} For an evaluator, the gain is a defensible bound to state to a lawyer in place of the case-by-case argument that suppression requires: each release spends some of a fixed \emph{privacy budget}, the parameters used to account for privacy loss across releases. The cost is that the added noise degrades exactly the small-area estimates a program audience most wants, so the method fits an open data portal that publishes cut after cut as well as other settings that require a formal privacy guarantee. Whichever standard a release uses, record it: a later reviewer needs to know whether the file shipped under a $k$ threshold or a privacy budget.

\section{Synthetic stand-ins: sharing structure, not records}\label{sec:synth}
When source records cannot be shared, a synthetic dataset\index{synthetic data} may let collaborators develop code inside a separate file. We build one here in a seeded call and inspect its fidelity and exact-match diagnostics. The source file remains in the approved environment, but the generated file still requires a disclosure assessment because it can reproduce source combinations. Use the synthetic file for development, run the final analysis on the source data in the approved environment, and label the two outputs clearly.\marginnote{An exact match is a review flag. Its risk depends on whether the source combination is rare, whether sensitive attributes also agree, what the recipient already knows, and the conditions of release.}

The box below introduces \texttt{synthgen}, which generates the stand-in and returns a fidelity and match diagnostic.

\begin{kaobox}[frametitle=Package Integration: \texttt{synthgen}]
\textbf{Install and run:} \texttt{synthgen}\index{synthgen@\texttt{synthgen}} generates that stand-in in one call, with no Python and no external dependencies. Its engine samples values from the observed marginal distributions and uses a joint normal draw to approximate rank associations, then introduces missing values. Check the resulting distributions and associations in the returned diagnostics; a simulated sample need not reproduce them exactly. It refuses string and ID-shaped variables. Install it with the house idiom of Appendix~\ref{app:toolkit} (repository \texttt{synthgen-stata-public}); the syntax is \texttt{synthgen varlist, frame(name)|saving(file) [n(\#) seed(\#)]}. Each run returns \texttt{r(maxdmean)}, \texttt{r(maxdrho)}, and \texttt{r(dupes)}. These scalars describe fidelity and exact matches; they do not estimate membership or attribute-disclosure risk.
\end{kaobox}

On the chapter's running stand-in, one seeded call synthesizes the four quasi-identifiers and the wage column together:
\par
{\footnotesize
\begin{verbatim}
sysuse nlsw88, clear
synthgen race married collgrad industry wage, ///
    frame(synth) seed(20260730)
\end{verbatim}
}
The diagnostic reports a worst standardized mean gap of 0.040, a worst rank-correlation gap of 0.103, and 421 of the 2{,}246 synthetic rows exactly matching a source row on these five columns. Part~(s) of \texttt{ch14\_kanon.do} compares those matches with source frequencies: 286 generated rows match a combination that occurs once in the source, and 409 match a combination that occurs no more than four times. These results support using this file only inside the development exercise. Do not release it externally without a case-specific disclosure assessment and additional protection. A uniqueness scan on the generated file alone cannot establish how rare a match is in the confidential source.

An outside team can use the synthetic file to develop code, then you can run that code against the confidential records in the approved environment. Before release, inspect exact matches and near matches, especially unusual combinations, and compare the planned analyses on real and synthetic data. A low overall match rate does not establish that the remaining matches are harmless. Describe the generator and the checks actually performed in the sharing documentation; do not promise that no synthetic record corresponds to a real person unless your assessment supports that statement.

\section{Sanitizing raw files with a human in the loop}\label{sec:rawsweep}
Protected information most often enters at intake, in the raw files a partner drops in the shared drive, so the gate belongs there. We suggest a human-in-the-loop sweep: scan each incoming file, emit a review workbook listing the fields the scan flags for removal, let a person confirm, then execute the removals with a dry-run mode and an audit receipt recording exactly what was stripped and when.\marginnote{Keep the human in the loop deliberately: an auto-redactor that silently drops a column is the fastest way to lose the one field you needed. The reviewer's confirmation is also what makes the receipt meaningful, since it records a decision, not just a script's guess.} The workflow is simple enough to script, and the box below shows the command that packages its scanning step.

\begin{kaobox}[frametitle=Package Integration: \texttt{rawsweep}]
\textbf{Integrated command:} \texttt{rawsweep}\index{rawsweep@\texttt{rawsweep}}. Files that land in a shared intake folder deserve a triage read before anything imports them. By default, \texttt{rawsweep} inventories CSV and TXT files and scans only their first 25 lines against the same six pattern classes as \texttt{ai\_privacy\_gate}. It can flag an early, recognizable pattern; it cannot establish that a file lacks protected information. Files in other formats, later lines, multiline fields, names, and indirect identifiers require separate review. The command does not delete or edit anything. On a three-file intake folder (part (r) of \texttt{ch14\_kanon.do}):
\par
{\footnotesize
\begin{verbatim}
. rawsweep, directory("raw/intake") pii clear

rawsweep: 3 file(s) in raw/intake
  smallest            0  largest           43 bytes
  1 zero-byte file(s): failed downloads wearing success flags
  files with likely PII in the first 25 rows: 1
  route those through the privacy gate before anything
    else touches them.
\end{verbatim}
}
The zero-byte flag is Chapter~\ref{ch:noapi}'s manifest triage, automated; the PII flag hands the file to Chapter~\ref{ch:ai}'s gate.
\end{kaobox}

The scan receipt records the files and lines inspected and the patterns flagged. Keep a separate record of the reviewer's decisions and the code that creates the cleaned copy. Run fuller inspection before the data enter derived files; the default triage result is one input to that review, not a release decision.

\subsection{The disclosure-review checklist: one page a lawyer and an evaluator can both read}\label{ch14-sc-checklist}
Before a file ships, someone has to sign, and the signer wants everything on one page. The disclosure-review checklist\index{disclosure-review checklist} is that page\marginnote{The checklist is this chapter's instance of the tracking spine of Chapter~\ref{ch:workshop} (Section~\ref{sec:ch02-spine}), with the columns renamed for release work.}: one row per release, six fields, in words neither side has to translate for the other.
\begin{itemize}
\item \emph{Quasi-identifiers listed:} the agreed columns in ordinary words, with the date and the names of who agreed (Section~\ref{sec:qid}).
\item \emph{$k$ computed:} the value \texttt{riskscan} returned, the file version it ran on, and the script that ran it.
\item \emph{Threshold applied:} the number and the document that set it, a data-use agreement clause or a partner policy, not an analyst's preference.
\item \emph{Suppression logged:} which cells were blanked, primary or complementary, and which exit each failing table took (coarsen, suppress, or aggregate) and why.
\item \emph{Sign-off named:} one person, by name and date; a team cannot sign.
\item \emph{Release archived:} the exact file that shipped, alongside the script and the intake receipt, so the release can be audited later.
\end{itemize}
You reach for the page at two moments. Before release, an incomplete row is reason enough to stop: a blank sign-off field or an unrecorded threshold is a hold, not a formality. A year later, when a partner asks what left the building and under what standard, the row is the answer.

\paragraph{Where this leaves you.} You can count quasi-identifier combinations, examine sensitive-value diversity, coarsen a proposed file, and apply primary and complementary suppression to a table. You can also generate a synthetic stand-in and screen intake files for likely protected information. Record the fields assessed, the release standard, the checks performed, and the reviewer responsible for the decision. These checks support disclosure review within their stated limits; overlapping releases and outside information still matter. Chapter~\ref{ch:tools} shows how to turn repeated do-file steps into shared commands and preserve the project in a replication archive.

\section*{Exercises}\addcontentsline{toc}{section}{Exercises}
\begin{enumerate}
\item \emph{Try it on your own data.} Take a de-identified file you plan to share, pick four quasi-identifiers a partner would plausibly hold, and build the cell size \texttt{k} the way the worked example does. Run \texttt{count if k==1} and report how many people in your file are one of a kind on those four columns.
\item Rerun \texttt{ch14\_kanon.do} on \texttt{nlsw88}, but swap the four quasi-identifiers for race, married, collgrad, and \texttt{occupation} in place of industry. Report the new count of unique records and say whether occupation is a finer or coarser mesh than industry.
\item Before you run anything, predict how many one-of-a-kind people remain if you coarsen industry to two groups (goods versus everything else) rather than three. Then compute it with \texttt{recode} and compare your guess to the actual count.
\item Break the complementary-suppression audit on purpose: take the agriculture column of the worked example, blank only the graduate cell, and leave the non-graduate cell visible against its total of 17. Confirm by hand that a reader recovers the hidden three by subtraction, then blank the second cell and confirm the leak closes.
\item Coarsen a single high-cardinality column in your own file until the $k=1$ bin is empty, and write one sentence naming exactly which analytic comparison the coarsening costs the reader and whether any planned table still needs that detail.
\item Install \texttt{riskscan} and run it on \texttt{nlsw88} with the four quasi-identifiers of the worked example, then read \texttt{r(below)} and write the one-line release gate (\texttt{if r(below) > 0}) you would put at the top of a publication do-file. Add \texttt{sensitive(union)} and report how many cells have $l=1$, then say in one sentence what an $l=1$ cell would leak that a large $k$ alone would not.
\item Draft the disclosure-review checklist row for the \texttt{nlsw88} release the worked examples build: fill all six fields (quasi-identifiers, $k$, threshold, suppression log, sign-off, archive path), and note which field you could not complete without a decision from someone other than the analyst.
\end{enumerate}
